\section{Demonstration of this template}

\subsection*{Hallo}
\subsection{Hallo1}

Given predictions $y_i$ and targets $t_i$, the mean squared error over $n$ samples is
\begin{equation}\label{eq:mse}
  \mathrm{MSE}=\frac{1}{n}\sum_{i=1}^n (y_i-t_i)^2.
\end{equation}
Some text yadayadayada.

\subsection{Custom boxes}
\subsubsection{Lemmas and proofs}

yadayadayada

yadayadayada

yadayadayada

yadayadayada

\begin{nicelemma}* \label{lem:whatever}
    Let $a,b,c\in\mathbb{Z}$. If $ca = cb$ and $c\neq 0$, then $a = b$.
    \nextpart[Example.]
    If $c > 0$, then $a = b$ is obvious. If $c < 0$, then $-c > 0$ and $-ca = -cb$, so $-a = -b$ and $a = b$.
    Just checking to see what happens if the proof ends in an equation
    \[
        \int_0^x y^2 \,\mathrm{d}x = \frac{1}{3} x^3 \,. 
    \]
    Can I say more?
    \[
    \text{More equations} = \pi \,. 
    \]

    \begin{equation}
    \begin{split}
      A &= B + C\\
        &= D 
    \end{split}\qedtag
    \end{equation}

    \begin{align}
      a &= b + c \\
        &= d \qedtag
    \end{align}
\end{nicelemma}

\begin{nicelemma} \label{lem:seemstowork}
    Let $a,b,c\in\mathbb{Z}$. If $ca = cb$ and $c\neq 0$, then $a = b$. \manualqed
\end{nicelemma}

Some text yadayadayada

\begin{nicetheorem}* \label{theo:whatever}
    Let $a,b,c\in\mathbb{Z}$. If $ca = cb$ and $c\neq 0$, then $a = b$.
    \nextpart
    If $c > 0$, then $a = b$ is obvious. If $c < 0$, then $-c > 0$ and $-ca = -cb$, so $-a = -b$ and $a = b$.
    Just checking to see what happens if the proof ends in an equation
    \[
        \int_0^x y^2 \,\mathrm{d}x = \frac{1}{3} x^3 \,. 
    \]
\end{nicetheorem}

\begin{nicetheorem*}* \label{theo:whatever2}
    Let $a,b,c\in\mathbb{Z}$. If $ca = cb$ and $c\neq 0$, then $a = b$.
    \nextpart
    If $c > 0$, then $a = b$ is obvious. If $c < 0$, then $-c > 0$ and $-ca = -cb$, so $-a = -b$ and $a = b$.
    Just checking to see what happens if the proof ends in an equation
    \[
        \int_0^x y^2 \,\mathrm{d}x = \frac{1}{3} x^3 \,. 
    \]
\end{nicetheorem*}

\begin{nicetheorem} \label{theo:seemstowork}
    Let $a,b,c\in\mathbb{Z}$. If $ca = cb$ and $c\neq 0$, then $a = b$.
\end{nicetheorem}

Some text yadayadayada

\begin{nicecorollary}* \label{cor:whatever}
    Let $a,b,c\in\mathbb{Z}$. If $ca = cb$ and $c\neq 0$, then $a = b$.
    \nextpart
    If $c > 0$, then $a = b$ is obvious. If $c < 0$, then $-c > 0$ and $-ca = -cb$, so $-a = -b$ and $a = b$.
    Just checking to see what happens if the proof ends in an equation
    \begin{equation}
        \int_0^x y^2 \,\mathrm{d}x = \frac{1}{3} x^3 \,. \qedtag
    \end{equation}
\end{nicecorollary}

\begin{nicecorollary} \label{cor:seemstowork}
    Let $a,b,c\in\mathbb{Z}$. If $ca = cb$ and $c\neq 0$, then $a = b$.
\end{nicecorollary}

Some text yadayadayada

\begin{nicedefinition}[Circle roundness]
    A circle is always round, omg right! ja ja ja ja nee nee nee nee
\end{nicedefinition}

\begin{nicealgorithm}*[Fibonacci]
    The algorithm to calculate fibonacci numbers.
    \nextpart[Algorithm.]
    Here the pseudocode:
    \vspace{0.5em}
    \hrule
    \vspace{0.5em}
    \begin{algorithmic}[1]
        \Function{Fibonacci}{$n$}
            \If{$n \leq 1$}
                \State \Return $n$
            \Else
                \State \Return \Call{Fibonacci}{$n-1$} + \Call{Fibonacci}{$n-2$}
            \EndIf
        \EndFunction
    \end{algorithmic}
\end{nicealgorithm}

\begin{nicecode}*
  Demo of fibonacci in python:
  \nextpart[Code:]*

  \begin{lstlisting}[language=Python]
def fib(n):
    if n <= 1:
        return n
    else:
        return fib(n-1) + fib(n-2)
  \end{lstlisting}
\end{nicecode}

\section*{NJA}

\texttt{Code fonts test} \\
Referencing is easy, just use \inlinecode{\texttt{\backslash Cref\{\}}}, for example u can reference to \Cref{lem:whatever}. Ofcourse u need to label it correctly using \inlinecode{\texttt{\backslash label\{\}}}.

\section*{Figures}

\begin{figure}
    \centering
    \includegraphics[width=\linewidth]{figures/test2.jpeg}
    \caption{Test Image}
    \label{fig:test_image}
\end{figure}

U can reference as follows: \Cref{fig:test_image} \\
\\
Test inlinebox: A piece of \inlinetext{text} by using \texttt{\backslash inlinetext\{$\cdots$\}}. \\
\\
Now try a piece of \inlinecode{code -> \#2 != \#1} by using \texttt{\backslash inlindecode\{$\cdots$\}}.

A default grey \inlinetext{note}
and a colored one \inlinetext{warning}[orange!30].

Code: \inlinecode{foo()->bar} and \inlinecode{alert!}[yellow!25].

\subsection{Custom Equation Boxes}

Starred \inlinecode{*}, as in \inlinecode{\textbackslash begin\{equationbox\}*} results in an unnumbered equationbox, while no star adds an equation number.

% Example 1: Standard (Numbered, Black Border, Transparent Fill)
\begin{equationbox}
    \tilde{x}_i = V_m s_i
\end{equationbox}

% Example 2: Starred (No Number, Black Border, Transparent Fill)
\begin{equationbox}*
    \tilde{x}_i = V_m s_i
\end{equationbox}

An equationbox takes $2$ parameters that are \inlinetext{optional}[green!25] and alter the \inlinetext{linecolor} and the\newline \inlinetext{fill color}. They can be used as follows: 
\begin{center}
    \inlinecode{\textbackslash begin\{equationbox\}[<linecolor>(Optional)][<fillcolor>(Optional)]}    
\end{center}


% Example 3: Custom Line Color (Red Line, Transparent Fill)
\begin{equationbox}[red]
    E = mc^2
\end{equationbox}

When you want multiple equations in one box that are aligned with respect to a "=" for example, use \inlinecode{\textbackslash begin\{aligned\}} and \inlinecode{\&=}. Just aligning multiple lines in a single box $\longrightarrow$ use \inlinecode{\textbackslash begin\{gathered\}}. You need so separate lines by ending a line with \inlinecode{\textbackslash\textbackslash}.

% Example 4: Custom Line AND Fill (Blue Line, Yellow Fill)
\begin{equationbox}[blue][yellow!20]
    \begin{gathered}
        \nabla \cdot \mathbf{E} = \frac{\rho}{\varepsilon_0} \\
        \text{JA ANOTHER ONE} \\
        yada = \sum_{i=1}^{\infty} y_a^d \times a 
    \end{gathered}
\end{equationbox}

\subsection{Custom left-bar box}

\begin{barbox}
This is some text with a default gray bar and white background.
\end{barbox}

\begin{barbox}[red]
This is some text with a red bar and white background.
\end{barbox}

\begin{barbox}[blue][yellow!20]
This one has a blue bar and a light yellow background.
\end{barbox}

\begin{barbox}[][green!10]
This one has the default gray bar and a light green background.
\end{barbox}

\subsection{Custom rounded text-box}

\begin{roundbox}
I am Sam
\end{roundbox}

\begin{roundbox}[yellow!20] % custom background, no border
Sam I am
\end{roundbox}

\begin{roundbox}[green!10][blue] % background + border color
I do not like green eggs and ham
\end{roundbox}

\newpage
\subsection{Custom environments}

\subsubsection{Environment with text on one side and image on other side}

\begin{sidebyimage}{figures/test2.jpeg}[0.5\textwidth][My Caption][fig:mylabel]
    Here is some text that describes the image. Because I did not specify 'l' or 'r', 
    the logic defaults to the right side. The text is vertically centered relative 
    to the image.
\end{sidebyimage}

\begin{sidebyimage}[l]{figures/test2.jpeg}[0.5\textwidth]
    Now the image is on the left. I omitted the caption and label here, 
    so the image just appears by itself. The text width is automatically 
    calculated to fill the rest of the page.
\end{sidebyimage}

\begin{wrapimage}[r]{figures/test2.jpeg}[0.5\textwidth][My Caption]
    This is a long paragraph of text. Because I am using 'wrapimage', 
    the text will start next to the image on the left. Once the text 
    exceeds the height of the image, it will naturally wrap underneath 
    it and span the full width of the page, just like in a newspaper 
    or magazine article. ha ha ha ha ha h a ha ah ah ah ah ah ah ah ah ah ah ah ah ah ah ah ah ah ah ah ah ah ah ah ah ah ah ah ah ah ah ah ah ah ah ah ah ah ah ah ah ah ah ah ah ah ah ah ah ah ah ah ah ah ah ah ah ah ah ah ah ah ah ah ah ah ah ah ah ah ah ah ah ah ah ah ah ah ah ah ah ah ah ah ah ah ah ah ah ah ah ah ah ah ah. ha ha ha ha ha h a ha ah ah ah ah ah ah ah ah ah ah ah ah ah ah ah ah ah ah ah ah ah ah ah ah ah ah ah ah ah ah ah ah ah ah ah ah ah ah ah ah ah ah ah ah ah ah ah ah ah ah ah ah ah ah ah ah ah ah ah ah ah ah ah ah ah ah ah ah ah ah ah ah ah ah ah ah ah ah ah ah ah ah ah ah ah ah ah ah ah ah ah ah ah ah. 
\end{wrapimage}